\thispagestyle{empty}
\vspace*{\fill}

\begin{flushleft}
{\sffamily\small

\textbf{Building Awesome MCP Apps}\\
\emph{What every AI application developer should know about the Model Context Protocol}

\vspace{12pt}
Copyright \textcopyright{} 2026 Krimler. All rights reserved.

\vspace{12pt}
Published online at\\
\texttt{https://github.com/krimler/Building-Awesome-MCP-Apps}

\vspace{12pt}
Written against Model Context Protocol revision \texttt{2026-07-28}.
Chapter~\ref{ch:architecture} explains the feature lifecycle policy that
governs how long that pinning stays useful, and what happens when it stops.

\vspace{12pt}
\textbf{The companion code.} Every listing in this book is extracted from the
Meridian reference application, which ships in the \texttt{meridian/} directory
of the repository above. It has no third-party runtime dependencies, its test
suite passes, and the numbers printed in these pages come from its measurement
harness rather than from memory. Appendix~\ref{ch:companion} is the tour.

\vspace{12pt}
\textbf{Cover.} Generated by \texttt{figures-src/plots/cover.py}, which is part
of this repository. It draws the traffic between five hosts and eleven servers,
which is the N${\times}$M problem Chapter~\ref{ch:primer} opens with. No
third-party artwork was used, so the provenance is unambiguous.

\vspace{12pt}
\textbf{Epigraphs.} Each chapter opens with a short line of film dialogue,
quoted for commentary. They are isolated in single macro calls, one per
chapter, so they can be replaced without touching the prose.

\vspace{12pt}
Typeset with \LaTeX{}. The colophon at the back has the details.

\vspace{12pt}
First edition.
}
\end{flushleft}

\cleardoublepage
